\section{Parameter asymptotics near the nonexceptional boundary}
\label{app:bipodal-parameter-expansions}

\begin{remark}[First-order asymptotics of the bipodal parameters]
\label{rmk:bipodal-parameter-expansions}
Part~\textup{(d)} of \Cref{thm:local-optimizer-structure} gives the qualitative limit
\(c(p,r)\to0\).  The purpose of this remark is to quantify that limit and the
corresponding change in the block densities.  Fix \(r\) in the neighbourhood
furnished by the theorem and let \(p\uparrow\pc(r)\) from the
symmetry-breaking side.  Relabel the podes as in \Cref{thm:local-optimizer-structure}, so
that the first pode has size \(c=c(p,r)\), and set
\[
  \delta_*:=r-e(W_{p,r}),
  \qquad
  z:=\zeta_d(r).
\]
Denote by \(q_{11}^0(r)\) the value at
\((\eps,\vartheta)=(r,0)\) of the analytically extended KRR--S parameter
\(q_{11}\) from \Cref{thm:krrs-analytic-extension}.
All limits and \(O\)-terms below refer to this regime and are uniform after
the \(r\)-neighbourhood is made smaller if necessary.

For \(z\ne r\), define
\begin{equation}\label{eq:edge-deficit-coefficient}
  D_d(r,z):=\frac{z^d-r^d}{d r^{d-1}}-(z-r)>0,
\end{equation}
where positivity follows from the strict convexity of \(u\mapsto u^d\).
Then
\[
  \delta_*=2D_d(r,z)c+O(\delta_*^2).
\]
Thus, \(2D_d(r,z)\) is the first-order edge-density deficit per unit of
smaller-pode mass, after \(q_{22}\) is adjusted to preserve the
\(H\)-density.  Combining this relation with
\(\delta_*=\lambda(p,r)/A_H(r)+O(\lambda(p,r)^2)\) from
\Cref{thm:local-optimizer-structure} gives
\begin{equation}\label{eq:optimizer-block-size}
  c(p,r)=
  \frac{\lambda(p,r)}
  {2D_d(r,\zeta_d(r))A_H(r)}
  +O(\lambda(p,r)^2),
\end{equation}
\begin{equation}\label{eq:optimizer-block-density}
  q_{22}(p,r)
  =
  r-
  \frac{\zeta_d(r)^d-r^d}
       {d r^{d-1}D_d(r,\zeta_d(r))A_H(r)}
  \lambda(p,r)
  +O(\lambda(p,r)^2),
\end{equation}
and
\[
  q_{12}(p,r)=\zeta_d(r)+O(\lambda(p,r)),
  \qquad
  q_{11}(p,r)=q_{11}^0(r)+O(\lambda(p,r)).
\]
In particular, both the smaller-pode size and the displacement of the
large-pode density from \(r\) are of order \(\lambda(p,r)\).
The other two block densities satisfy
\(q_{12}\to\zeta_d(r)\ne r\) and
\(q_{11}\to q_{11}^0(r)\).  These relations quantify the vanishing-pode
mechanism described in the introduction.

The limiting cross density \(\zeta_d(r)\) should not be confused with the
Maxwell contact density \(\sm(r)\) used in the comparison graphons of
\Cref{sec:nondegeneracy}.  The former is selected by the KRR--S entropy
maximization, whereas the latter is determined by the Maxwell contact
equations \eqref{eq:contact-equal-slopes}--\eqref{eq:contact-chord-identity}.  They arise from
different variational problems and need not agree.

For completeness, we derive the two explicit expansions above.  Suppress the
arguments \((p,r)\) and evaluate the analytic KRR--S parameter map at
\[
  (\eps,\vartheta)
  =\bigl(r-\delta_*,\,r^m-(r-\delta_*)^m\bigr).
\]
Since \(\vartheta=O(\delta_*)\), \Cref{thm:krrs-analytic-extension} and analyticity at
\((r,0)\) give
\[
  c=O(\delta_*),\qquad
  q_{22}-r=O(\delta_*),\qquad
  q_{12}-z=O(\delta_*),\qquad
  q_{11}-q_{11}^0(r)=O(\delta_*).
\]
The exact edge-density identity
\[
  e(W_{p,r})
  =q_{22}+2c(q_{12}-q_{22})
   +c^2(q_{11}-2q_{12}+q_{22})
\]
therefore yields
\begin{equation}\label{eq:bipodal-edge-linearization}
  -\delta_*=(q_{22}-r)+2(z-r)c+O(\delta_*^2).
\end{equation}

For the \(H\)-density, separate the homomorphisms according to the number of
vertices mapped to the smaller pode.  Since \(H\) is \(d\)-regular and has
\(2m/d\) vertices, the contributions with zero and one such vertex give
\begin{align*}
  t(H,W_{p,r})
  &=(1-c)^{2m/d}q_{22}^m
    +\frac{2m}{d}c(1-c)^{2m/d-1}
      q_{12}^d q_{22}^{m-d}+O(c^2)\\
  &=r^m+m r^{m-1}(q_{22}-r)
    +\frac{2m}{d}r^{m-d}(z^d-r^d)c
    +O(\delta_*^2).
\end{align*}
Using \(t(H,W_{p,r})=r^m\), we obtain
\begin{equation}\label{eq:bipodal-density-linearization}
  0=m r^{m-1}(q_{22}-r)
    +\frac{2m}{d}r^{m-d}(z^d-r^d)c
    +O(\delta_*^2).
\end{equation}
Eliminating \(q_{22}-r\) between
\eqref{eq:bipodal-edge-linearization} and
\eqref{eq:bipodal-density-linearization} gives
\[
  \delta_*=2D_d(r,z)c+O(\delta_*^2),
  \qquad
  c=\frac{\delta_*}{2D_d(r,z)}+O(\delta_*^2),
\]
which proves \eqref{eq:optimizer-block-size}.  Finally,
\eqref{eq:bipodal-edge-linearization} gives
\[
  q_{22}
  =r-\frac{z^d-r^d}{d r^{d-1}D_d(r,z)}\delta_*
   +O(\delta_*^2),
\]
and \eqref{eq:optimizer-block-density} follows.
\end{remark}